\section{Discussion}
\label{sec:discussion}

\subsection{Why This Waste Exists}

The append-only pattern in agentic AI tools is not an implementation
bug. It arises from four reinforcing factors:

\begin{enumerate}[nosep]
  \item \textbf{Training data}: Models learn from conversations that
    grow monotonically. No training examples demonstrate intelligent
    content removal.
  \item \textbf{API design}: The Messages API accepts a list of
    messages. The natural operation on a list is append. The API
    provides no mechanism for ``this content is stale.''
  \item \textbf{Framework defaults}: Every orchestration framework
    appends by default. Eviction requires explicit engineering.
  \item \textbf{Invisible cost}: Token consumption is billed after
    the fact. Quality degradation from context bloat is invisible
    unless measured. There is no backpressure signal.
\end{enumerate}

\subsection{The Inverted Cost Model}
\label{sec:inverted-cost}

In traditional virtual memory, keeping a page in physical memory is
free. RAM is allocated whether the page is accessed or not; the only
cost is \emph{faulting}---the latency of disk I/O when a needed page
is not resident. The entire replacement algorithm literature, from
FIFO to LRU to Clock to Working Set, optimizes a single objective:
\emph{minimize faults}.

In LLM context management, the cost model is inverted. Every token
kept in context costs money on every turn. A 5{,}000-token file
sitting in context for 20~turns costs 100{,}000 input tokens of
processing. Re-reading that file once when actually needed costs
5{,}000 tokens. The eviction saves 95{,}000 tokens. Keeping is
expensive; faulting is cheap.

This inversion changes the optimization objective. The correct
formulation is not ``minimize faults'' but:
\[
  \min \sum_{p \in P} \bigl[ C_{\text{keep}}(p) + C_{\text{fault}}(p) \bigr]
\]
where $C_{\text{keep}}(p) = |p| \cdot T_{\text{resident}}(p) \cdot
c_{\text{token}}$ is the cumulative cost of keeping page $p$ resident
for $T_{\text{resident}}$ turns, and $C_{\text{fault}}(p) = |p| \cdot
c_{\text{token}}$ is the one-time cost of restoring the page (one
turn of reprocessing). The break-even condition is:
\[
  |p| \cdot T_{\text{until\_next\_ref}}(p) \cdot c_{\text{token}} >
  |p| \cdot c_{\text{token}}
\]
which simplifies to: \emph{evict whenever the page will not be
referenced for more than one turn}. This is why FIFO works so well
in our system despite being the worst-performing policy in classical
VM: when keeping is expensive and faulting is cheap, aggressive
eviction is correct by default. Sophistication in replacement policy
is needed only to avoid evicting pages that will be needed on the
\emph{very next} turn.

\paragraph{Belady's MIN under inverted costs.}
Belady's optimal algorithm~\citep{belady1966study} evicts the page
whose next reference is farthest in the future, minimizing total
faults. Under the inverted cost model, the optimal policy is
different: it must minimize the sum of keeping costs and fault costs.
A page with a distant next reference should be evicted immediately
(saving many turns of keep cost at the price of one fault). A page
referenced every turn should never be evicted (the fault cost equals
one turn of keep cost, so eviction saves nothing). The optimal
offline policy under inverted costs is not MIN---it is a
cost-weighted variant where pages are evicted when their projected
keep cost exceeds their fault cost, regardless of whether a fault
will occur.

This has a practical implication: the \emph{size} of the page matters
for eviction priority in a way it does not in classical VM (where all
pages are equal size). A 10{,}000-token file costs ten times as much
per turn as a 1{,}000-token file. Large pages should be evicted
eagerly unless access frequency justifies the keep cost.

\paragraph{Non-linear fault cost.}
The linear fault cost $C_{\text{fault}}(p) = |p| \cdot c_{\text{token}}$
above is a simplification. Transformer inference is dominated by
self-attention, which is $O(n^2)$ in sequence length. A page fault
requires an additional full inference pass---the model emits a
\texttt{tool\_use} for the recall handle, and the restored content
returns as a \texttt{tool\_result} message that triggers a second
inference over the entire context. The true fault cost is therefore
proportional to $n^2$ at the current context size $n$, not to the
page size $|p|$ alone.

This produces a counter-intuitive policy gradient. At low fill
(e.g., $n = 40\text{K}$), faults are cheap and aggressive eviction is
correct---the quadratic cost of an extra pass is modest, and keeping
the working set small reduces the $n^2$ base for all subsequent turns.
At high fill (e.g., $n > 100\text{K}$), faults become expensive: the
additional inference pass costs roughly $(n + |p|)^2 \approx n^2$
tokens of compute. The eviction policy should therefore become
\emph{more conservative} as context pressure rises---the opposite of
the naive instinct to evict aggressively under pressure. At high
fill, the system should evict only content it is highly confident will
not be referenced again.

\paragraph{Object-based cost variance.}
Unlike classical VM, where all pages are fixed-size blocks (typically
4\,KB), Pichay manages variable-size objects: a recalled tensor may be
a 3-line summary or a 200-line file. This means eviction and fault
costs vary by orders of magnitude across the working set. The
replacement policy cannot treat all eviction candidates as equivalent;
it must weight both the per-turn keep cost and the potential fault cost
by object size. Combined with the quadratic fault penalty, this argues
for a \emph{size-aware, fill-sensitive} replacement policy that has no
direct analogue in the classical VM literature.

\paragraph{Cache invalidation cost.}
The cost model so far considers only token-level costs (keeping and
faulting). Structural mutations---collapse operations that remove or
replace blocks in the message array---have an additional cost:
\emph{prompt cache invalidation}. Inference providers cache the
tokenized prefix of repeated requests; when the prefix changes, the
cache misses and the entire context must be reprocessed.

In production, a collapse operation that compressed 12~turns of
orientation dialogue into a single summary sentence caused the cache
hit rate to drop from 100\% to 25\% for one turn, then recover to
100\% on the following turn as the new prefix stabilized. The cost
was one full recompute of $\sim$105K tokens---comparable to the
cost of several page faults. A collapse that saves 10KB of context
but invalidates a 100K-token cached prefix is a net loss unless the
space savings persist for enough subsequent turns to amortize the
one-time recompute.

This argues for \emph{batching} structural mutations: accumulate
collapse candidates and execute them together in a single pass,
paying the cache invalidation cost once rather than per-operation.
It also argues against frequent small collapses in favor of
infrequent large ones.

\paragraph{Pin decay.}
The inverted cost model also argues against permanent pins. In our
current design (Section~\ref{sec:pinning}), one fault pins a page
permanently. But a fault tells us the content was needed \emph{then},
not forever. Under the cost model, a pin should decay: its strength
halves every $K$~turns since last access, and the page becomes
evictable when the projected keep cost of the remaining pin lifetime
exceeds the fault cost. This gives LRU-like behavior with
cost-weighted decay, and prevents the monotonic working-set growth
that permanent pinning causes in long sessions.

\subsection{Generalizability}

The waste patterns we measure are structural consequences of the
agentic architecture:

\begin{itemize}[nosep]
  \item Any system that sends tool definitions on every request
    will send schemas for unused tools. The specific percentage
    depends on tool count and usage distribution, but the mechanism
    is universal.
  \item Any system that injects instructions into messages risks
    duplication. Claude Code's triplication is egregious but not
    unique.
  \item Any system that keeps tool results in context for the
    session lifetime will see amplification proportional to session
    length.
\end{itemize}

Our empirical evidence is from one system (Claude Code) and one
user. The structural argument is strong---the architecture
forces the waste---but comparative measurement of other tools
would strengthen the generalizability claim.

\subsection{Scope, Non-Goals, and Threats to Validity}

\paragraph{Scope.}
This paper evaluates context waste and demand paging behavior for
agentic coding workloads with persistent tool and transcript history.
The target claim is mechanism-level: when requests repeatedly resend
large static prefixes and stale tool output, eviction plus fault-based
restoration reduces cumulative input processing.

\paragraph{Non-goals.}
We do not claim universal percentage reductions across all assistants,
all model families, or all task domains. We do not evaluate frontier
reasoning benchmarks, autonomous planning quality, or human preference
at product scale. We also do not claim that L1 expansion (larger
context windows) is obsolete; we claim it is incomplete without
hierarchical management.

\paragraph{Threats to validity.}
External validity is limited by the single-user corpus and one primary
assistant implementation. Internal validity risks include prompt drift,
toolset changes across versions, and workload phase effects that alter
hot-set composition. To mitigate drift, we report a frozen cohort for
headline numbers and provide executable recount tooling plus
timestamped live snapshot artifacts for auditability.

\subsection{Quality: LLM-Judged Equivalence Check}
\label{sec:quality}

The intervention must not degrade output quality. We therefore run a
paired equivalence-style check: if results are broadly comparable while
compute drops, the intervention is useful.

\paragraph{Protocol.}
We select 18~sessions from the corpus (426--1{,}354~messages,
340K--718K~characters of context). For each session, we construct
paired contexts at 65--75\% of the conversation: the \emph{baseline}
retains all messages; the \emph{treatment} replaces consumed tool
results outside a 20-message recency window with tombstones, modeling
\textsc{Pichay}'s eviction behavior. Both conditions receive the same
continuation prompt (the next user message). Outputs are generated by
Sonnet~4 under both conditions.

An ensemble of three LLM judges (one Sonnet~4, two Haiku~4.5)
evaluates each paired output on correctness, completeness, and
coherence (1--5~scale). Judges also state a blind preference
(A, B, or tie) and whether they can identify which output had
reduced context. A/B assignment is randomized per judge to
prevent position bias; scores are remapped to baseline/treatment
before aggregation.

\paragraph{Results.}

\begin{table}[ht]
\centering
\caption{Non-inferiority evaluation: 18~sessions, 54~verdicts.
  Treatment evicts consumed tool results outside a 20-message
  recency window (mean compression 48\%).}
\label{tab:noninferiority}
\small
\begin{tabular}{@{}lrr@{}}
\toprule
\textbf{Metric} & \textbf{Baseline} & \textbf{Treatment} \\
\midrule
Judges preferring & 15 (28\%) & 20 (37\%) \\
Ties & \multicolumn{2}{c}{19 (35\%)} \\
\midrule
Mean correctness & 3.89 & 3.74 \\
Mean completeness & 3.59 & 3.59 \\
Mean coherence & 3.74 & 3.69 \\
Max score $\Delta$ & \multicolumn{2}{c}{0.15} \\
\midrule
Detection rate & \multicolumn{2}{c}{57\% ($p = 0.14$, n.s.)} \\
\bottomrule
\end{tabular}
\end{table}

Judges prefer the treatment \emph{more often} than the baseline
(37\% vs.\ 28\%), with 35\% ties. Completeness scores are identical;
correctness and coherence differ by at most 0.15~points on a 5-point
scale. Detection rate (57\%) is not significantly above chance
(binomial $z = 1.09$, $p = 0.14$, $n = 54$).

\paragraph{Failure modes.}
Two sessions (11\%) produced degenerate treatment outputs (zero
characters): the continuation prompt implicitly referenced content
that existed only in a tombstoned tool result. These are genuine
eviction casualties. The failure pattern is specific: the user's
prompt references a recently-consumed tool result by content rather
than by name. A reference-aware eviction policy would detect such
forward dependencies; the current heuristic (consumed + outside
recency window = evictable) does not.

\paragraph{The attention-concentration effect.}
In several sessions, all three judges preferred the treatment output.
Removing half the context did not merely preserve quality---it
improved it. This is consistent with the attention-dilution hypothesis:
transformer attention distributes weight across all tokens in context;
irrelevant tokens (stale tool results, consumed intermediate outputs)
dilute attention on the tokens that matter. Evicting them concentrates
attention on signal. The effect is strongest in sessions with high
tool-result density, where the evicted content is genuinely noise.
We note this as preliminary evidence; establishing the effect
rigorously would require controlled attention analysis with larger~$N$.

\subsection{Compute and Energy Implications}
\label{sec:energy}

The 970~million tokens of addressable waste in our corpus have
three compute cost dimensions:

\paragraph{Input processing.}
At a 93.5\% cache hit rate, most saved tokens are cache reads
(907M tokens at reduced cost). However, 63M tokens are cache
creation or uncached---full forward-pass computation that need not
occur.

\paragraph{Attention cost.}
For every output token, the model computes attention over the full
context. Removing 17{,}913 tokens from a mean context of 82{,}061
reduces the per-output-token attention computation by 21.8\%.
Across the corpus, this eliminates 85~billion token-token attention
pairs.

\paragraph{Throughput.}
Each token in context requires KV cache memory. For large models,
the per-token KV cache footprint is on the order of megabytes.
Smaller context per request means less KV cache memory per request,
which means more concurrent requests per GPU. A 21.8\% context
reduction translates, to first approximation, to 21.8\% more
concurrent users on the same GPU fleet.

This is the argument that changes the infrastructure economics.
The savings are not just ``less cost per request''---they are
``more requests per machine,'' which is the difference between
building new data centers and not building them.

\paragraph{Compounding across sessions.}
These per-call savings compound. Waste prevented on turn~$N$ is not
merely saved once---it is absent from every subsequent turn's
context, where it would have been reprocessed under $O(n^2)$
attention cost. The total compute savings across a session are
therefore superlinear in the waste fraction and grow with session
length. For a 100-turn session with 21.8\% addressable waste, the
cumulative attention savings are not 21.8\% of total session
compute---they are substantially larger, because the quadratic cost
of carrying unnecessary tokens compounds on every subsequent call.
The infrastructure implication follows directly: the fastest tokens
are the ones you never process.

Figure~\ref{fig:cumulative-cost} illustrates this compounding on an
88-turn live session running through \textsc{Pichay}. Cumulative
baseline cost (processing all context every turn) reaches 8.6M~input
tokens; managed cost reaches 4.8M---a 45\% cumulative reduction,
substantially larger than the per-turn compression ratio, because
each evicted token is absent from \emph{every subsequent turn}.

\begin{figure}[t]
\centering
\begin{tikzpicture}
\begin{axis}[
    width=\columnwidth,
    height=0.6\columnwidth,
    xlabel={Turn},
    ylabel={Cumulative input tokens (millions)},
    xmin=0, xmax=88,
    ymin=0, ymax=9.5,
    legend pos=north west,
    legend style={font=\small},
    grid=major,
    grid style={gray!20},
    scaled y ticks=false,
    y tick label style={/pgf/number format/fixed,
                        /pgf/number format/precision=1},
    yticklabel={\pgfmathparse{\tick/1000000}\pgfmathprintnumber{\pgfmathresult}},
]

% Baseline: cumulative tokens without eviction
\addplot[thick, red!70!black, mark=none, name path=baseline] coordinates {
    (0, 13651) (10, 352856) (20, 1037825) (30, 1850237)
    (40, 2884656) (50, 3989325) (60, 5146048) (70, 6379942)
    (80, 7689415) (87, 8645027)
};
\addlegendentry{Baseline (full context)}

% Managed: cumulative tokens with Pichay eviction
\addplot[thick, blue!70!black, mark=none, dashed, name path=managed] coordinates {
    (0, 7508) (10, 194065) (20, 570793) (30, 1017614)
    (40, 1586539) (50, 2194101) (60, 2830295) (70, 3508933)
    (80, 4229138) (87, 4754720)
};
\addlegendentry{Managed (Pichay)}

% Shade the area between curves
\addplot[red!10, opacity=0.5, forget plot]
    fill between[of=baseline and managed];

\end{axis}
\end{tikzpicture}
\caption{Cumulative input tokens processed over an 88-turn session.
  Without context management, every turn re-processes all prior
  content, producing superlinear cumulative cost. Pichay's eviction
  reduces cumulative token spend by 45\% (\$11.67 at Opus pricing).
  The shaded area represents wasted computation.}
\label{fig:cumulative-cost}
\end{figure}


\paragraph{Fleet-scale extrapolation.}
Our corpus represents one user over four months. A fleet-wide
projection requires per-user token consumption data that inference
providers hold but have not published. The extrapolation framework
is straightforward: $N$ users $\times$ 970M tokens per user $=$
total addressable waste. The mechanism (tool schema waste, content
duplication, result amplification) is architectural, so the
per-user waste fraction should be relatively stable even if
absolute token counts vary.

% ====================================================================
